\scanpage{209}\mspage{154}
Si les deux étaient \deleted{\unsure{égaux}} \inserted{iso dans $\mathrm{HOT}_{\deleted{\unsure{S}}\unsure{W|w}}(S)$} par la flèche canonique
\[
 X_e\longrightarrow X,
\]
en appliquant le foncteur $p_!$ (où $p:S\longrightarrow e$) il s'ensuivrait que $X_e\longrightarrow X$ est dans $W$. Si par exemple $S=\unsure{\Delta|\mathbb{\Delta}}^1$, cela impliquerait que pour tout $X$ lisse sur $S$, $X_1\longrightarrow X$ est \deleted{dans} $W$, alors que visiblement la lissité implique que $X_0\longrightarrow X$ est dans $W$. \deleted{Si la lissité implique} Donc ça reviendrait à dire que \deleted{$X$ lisse} $\underline{\mathrm{Liss}}\,\Delta^1=\deleted{\unsure{\mathrm{Fib}\Delta^1}}\underline{\mathrm{Liss}}_0\unsure{\{\Delta^1\}|\Delta^1}$, d'où aussitôt $\underline{\mathrm{Liss}}\,S=\underline{\mathrm{Liss}}_0\,S$ \deleted{\unsure{ça se}} pour \deleted{tout} $S$ dans $\mathrm{Cat}$, ça se saurait ! Le contre \unsure{exemple} le plus trivial est de prendre $X=\{0\}\lhook\joinrel\longrightarrow S=\Delta^1$, l'inclusion de l'élément ``générique'', immersion ouverte, donc lisse, la fibre \deleted{\unsure{en}} $X_1$ est vide !

\margin{$\unsure{\mathrm{Z}|2}$}
\uline{Corollaire 1 rectifié} Si $X$ est dans $\mathrm{Cat}_{/S}$, pas néc. lisse sur $S$, alors le foncteur
\[
 \mathrm{mult}_{X/S}:Y\unsure{\longrightarrow|\longmapsto} Y\times_S X:\deleted{\unsure{\mathrm{Cat}}}\mathrm{Cat}_{/S}\longrightarrow\mathrm{Cat}_{/S}
\]
induit un foncteur
\[
 \underline{\mathrm{Liss}}\,S\longrightarrow\mathrm{Cat}_{/S}
\]
compatible avec les localiseurs $W_S^d,W_S^d$, donc induit un foncteur
\[
 (\underline{\mathrm{Liss}}\,S)(W_S^d)^{-1}\longrightarrow(\mathrm{Cat}_{/S})(W_S^d)^{-1}
\]
% archive sheet number: 208
